% Transcription — fonds Alexandre Grothendieck, shelfmark 51, batch 2 (pages 21-22).
% Established from the facsimile archives/batches/51/batch-02.pdf (a mirror of
% https://grothendieck.umontpellier.fr/51.pdf), per the skill
% transcribe-grothendieck. The folder is twenty-two pages; this batch is its
% last two. The batch file carries no cover sheet: PDF page N is page 20+N of
% the folder (the Montpellier footer prints « 21 » and « 22 »), and the
% archivists' pencil numbers « 21 » and « 22 », bottom left, agree.
%
% Pass: Opus 5.5 (claude-opus-5-5), 2026-09-23 — first pass, unchecked against the pages by a human.
%
% Both pages are transcribed; nothing is absent. Page 21 carries his own
% number « 11 » top right, so it continues a run he paginated from batch 1
% (recorded once in a \note). Page 22 is a fresh sheet under a heading of his
% own, « Schémas abéliens sur un S régulier », which becomes the second
% \section; the first \section heading is bracketed, being ours. The folder
% title « " Bonne réduction " (Shimura-Koizumi) » is not bracketed in the
% inventory and copies the words on the folder's cover sheet (batch 1, p. 1).
%
% What the batch is.
%   21  End of the argument begun in batch 1 (p. 20 closes on the question
%       whether u : A -> P, P the image in Pic_{X/S}, is a closed immersion):
%       « Donc l'image de A est P »; it remains to show that the kernel K of
%       u, finite over S, is the unit group. The special fibre K_0 is a
%       p-group, p the characteristic exponent of the residue field; a boxed
%       N.B., cancelled by two long diagonals, argues it through a section of
%       order n prime to p lifting to S; then the clean version: for n prime
%       to p, (Pic_{X/S})_n is unramified over S and A_n is an étale cover of
%       S, so in residue characteristic 0 « on a fini »; otherwise « il faut
%       regarder de plus près ». Below a rule, a programme: kernels and images
%       of homomorphisms of abelian and « multiabéliens » schemes in
%       characteristic ≠ 0, notably over a non-reduced base.
%   22  « Schémas abéliens sur un S régulier »: (a) an abelian scheme over S
%       is projective (« Raisonnement de Nakai-… »); (b) a torsor P has an
%       abelian scheme G « de translations », and P is an isotrivial principal
%       homogeneous space under G; (c) H^1(S,G) ⊂ H^1(K,G), P_K extends iff it
%       « se réduit bien », and there is a largest dense open U over which it
%       extends; (d) S − U has codimension 1 (margin: « résulte du th. de
%       pureté »), whence H^1(K,G)/H^1(S,G) embeds in a sum over prime divisors
%       D of local terms, and a question in a loop: an analogue of the purity
%       theorem for good reduction of abelian schemes?
%
% The hand is fast drafting on both sheets; statements and formulas come
% through, much of the connective prose on p. 21 does not, and is left \ill{}.
% Notation fixed for the batch, as in batch 1: his underlined Pic is set
% \mathrm{Pic} without the underlining; K_0 is the special fibre of the
% kernel; words he underlines in prose are set in \emph.
\documentclass[11pt,a4paper]{article}
% Le préambule commun à toutes les transcriptions du fonds Grothendieck.
%
% Il tient en une idée : **l'incertitude se compose, elle ne se résout pas.**
% Une transcription qui lisse un mot douteux a détruit la seule chose pour
% laquelle on la faisait — un lecteur doit pouvoir savoir, à chaque endroit,
% ce qui était sur le feuillet et ce qui a été ajouté. Les six macros qui
% suivent sont donc l'appareil critique, et non de la décoration.
%
% Les mêmes commandes sont reconnues par `scripts/render.mjs`, qui produit la
% vue de lecture du site. Une macro ajoutée ici doit l'être aussi là-bas,
% faute de quoi le rendu échouera bruyamment — ce qui est voulu : un rendu
% silencieusement faux est pire qu'un rendu absent.

\ProvidesPackage{grothendieck}[2026/08/08 Transcriptions du fonds Grothendieck]

\usepackage{amsmath,amssymb,amsthm}
\usepackage{mathrsfs}
\usepackage{stmaryrd}
% Les diagrammes commutatifs sont partout dans ce fonds : c'est la langue
% même de la géométrie algébrique de Grothendieck, et une transcription qui
% les rendrait en prose ne transcrirait rien.
\usepackage{tikz-cd}
% Français par défaut — c'est la langue des feuillets — mais l'anglais est
% chargé aussi : la traduction et le résumé sont des documents anglais, et la
% typographie française (espaces avant les deux-points) y serait une faute.
% Ces documents basculent par \selectlanguage{english} en tête de corps.
\usepackage[english,french]{babel}
\usepackage{fontspec}
\usepackage{geometry}
\geometry{a4paper,margin=25mm,marginparwidth=28mm}
\usepackage{marginnote}
\usepackage{xcolor}
% ulem, not soul. `soul`'s \st and \ul are built on a character-by-character
% scan that XeLaTeX's font machinery breaks: the document fails with an
% undefined control sequence rather than degrading. ulem does the same two jobs
% and is Unicode-engine safe. `normalem` stops it hijacking \emph, which the
% transcriptions use for Grothendieck's own emphasis.
\usepackage[normalem]{ulem}
\usepackage{hyperref}
\hypersetup{colorlinks=true,linkcolor=black,urlcolor=black,pdfborder={0 0 0}}

% --- Métadonnées du lot -----------------------------------------------------
% Reprises telles quelles par le script de rendu, qui les affiche en tête de
% la vue de lecture. Elles fixent ce que le fichier prétend couvrir : sans
% elles, une transcription flotte sans point d'attache dans un fonds de seize
% mille feuillets.

\newcommand{\folder}[1]{\gdef\grothFolder{#1}}
\newcommand{\batch}[1]{\gdef\grothBatch{#1}}
\newcommand{\leaves}[2]{\gdef\grothFirst{#1}\gdef\grothLast{#2}}
\newcommand{\foldertitle}[1]{\gdef\grothFoldertitle{#1}}
\gdef\grothFolder{?}\gdef\grothBatch{?}\gdef\grothFirst{?}\gdef\grothLast{?}\gdef\grothFoldertitle{}

% --- Le résumé --------------------------------------------------------------
%
% Il ouvre la lecture modernisée, sous le titre « Résumé », et il est composé
% en petit. Ce n'est pas un ornement : c'est le seul passage du document qui
% ne soit pas des mathématiques, et un lecteur doit pouvoir le reconnaître
% avant de l'avoir lu — puis le sauter s'il connaît déjà le sujet, ou s'y
% arrêter s'il ne le connaît pas. Le filet qui le suit marque la reprise du
% corps.

\newenvironment{resume}
  {\small\setlength{\parskip}{0.45em}}
  {\par\medskip\hrule\medskip}

% --- Mots-clés ---------------------------------------------------------------
%
% En anglais, en clôture du résumé de la lecture modernisée : le vocabulaire
% moderne sous lequel chercher ce que les feuillets construisent. Le manifeste
% du site (`npm run manifest`) les extrait pour étiqueter la cote entière —
% cette ligne est la source unique des tags, il n'y a pas de fichier à côté.
\newcommand{\keywords}[1]{\par\smallskip\noindent
  {\footnotesize\textsc{Keywords} --- \textit{#1}}\par}

% --- L'appareil critique ----------------------------------------------------

% Le numéro de feuillet, en marge. C'est l'ancre qui permet de régler un
% désaccord sur une formule en regardant *un* feuillet plutôt qu'une chemise
% de six cents. Le site s'en sert aussi pour tourner les pages du fac-similé
% quand on fait défiler la transcription.
\newcommand{\leaf}[1]{%
  \marginnote{\footnotesize\color{gray!70}\textsf{#1}}%
  \label{leaf:#1}\ignorespaces}

% Illisible : on ne devine pas. Un mot inventé qui se lit comme les autres est
% le pire résultat possible.
\newcommand{\ill}{\textcolor{red!60!black}{[\,\ldots\,]}}

% Lecture douteuse : le mot est donné, et signalé comme incertain.
\newcommand{\uncertain}[1]{\uline{#1}}

% Ajout de l'éditeur — une lettre manquante, un mot sous-entendu.
\newcommand{\add}[1]{\textcolor{blue!50!black}{[#1]}}

% Biffé par Grothendieck lui-même. Ce qu'il a rayé fait partie du document :
% une rature dit souvent où la pensée a changé de direction.
\newcommand{\struck}[1]{\sout{#1}}

% Note du transcripteur, sur l'état matériel du feuillet.
\newcommand{\note}[1]{\par\smallskip{\small\color{gray!60!black}\textit{[#1]}}\par\smallskip}

% Note marginale de Grothendieck, distincte du corps du texte.
\newcommand{\marginal}[1]{\par\smallskip\noindent\hspace{1em}%
  {\small\color{gray!60!black}#1}\par\smallskip}

% --- Titre ------------------------------------------------------------------

\AtBeginDocument{%
  \begin{center}
    {\large\bfseries Fonds Alexandre Grothendieck --- cote n\textsuperscript{o}~\grothFolder}\\[2pt]
    {\itshape\grothFoldertitle}\\[4pt]
    {\small feuillets \grothFirst--\grothLast\ (lot \grothBatch)}\\[2pt]
    {\footnotesize\color{gray!60!black}Transcription établie d'après le fac-similé numérisé
     par l'Université de Montpellier.\\
     Les lectures incertaines sont soulignées, les illisibles notées [\,\ldots\,],
     les ajouts entre crochets.}
  \end{center}
  \vspace{1em}}

\endinput


\folder{51}
\batch{2}
\pages{21}{22}
\dating{s.d. — le groupe « Variétés abéliennes » (45 à 56) est daté 1961-1973}
\watermark{Édition de démonstration}
\foldertitle{« Bonne réduction » (Shimura-Koizumi) : notes manuscrites (s.d.).}

\begin{document}

\section{[Le noyau de $u : A \to P$]}
\note{titre de l'éditeur, entre crochets : le feuillet 21 n'en porte aucun et
poursuit l'argument du feuillet 20 (lot 1), qui s'achevait sur la question de
savoir si $A \to P$ est une immersion fermée d'image $P$.}

\page{21}
\note{pagination de l'auteur : « 11 », en haut à droite.}

Donc l'image de $A$ est $P$. Reste à vérifier que le noyau de
\struck{$A$ est} $u : A \to P$ \struck{$\to$} [qui, on le sait, est fini sur
$S$, \uncertain{plat} \uncertain{génériquement} \ill{} le \uncertain{groupe
unité}] est le groupe unité sur $S$.
\note{la fin de la parenthèse est très rapide ; « le groupe unité » y est lu
sur le modèle de la même expression, mieux formée, qui suit le crochet.}

En tous cas, \struck{ce groupe est d'ordre} \struck{\uncertain{Ce fib.}} si
$p$ est l'exposant caractéristique du corps résiduel $k$, alors la fibre $K_0$
est un $p$-groupe.

\struck{En effet, si $\xi_0$ est un pt de $K_0$ \ill{} \ill{}}
\note{ligne biffée d'un trait ; sous elle, une addition interlinéaire
reliée par un trait, « de \ill{} », ne se lit pas.}

[N.B. \ill{} \ill{}, \ill{} la \ill{} \uncertain{$\xi_0$} d'ordre $n$
premier à $p$, alors $\xi_0$ provient d'une section $\xi$ \uncertain{d'ordre
$n$} de $A$ sur $S$, qui donne une section \uncertain{d'ordre $n$} de
$\mathrm{Pic}_{X/S}$ \ill{} nulle \uncertain{au} pt
\uncertain{fermé} $s$, \uncertain{donc}\,\ill{}]
\note{tout ce N.B. est encadré et annulé par deux longues diagonales ; on le
transcrit sans le biffer mot à mot. Sous « qui donne une » s'insère une petite
addition, « \& $X$ \ill{} $/S$ », dont la place dans la phrase n'est pas
sûre.}

\uncertain{Ceci} résulte facilement du fait que si $n$ est premier à $p$,
alors $(\mathrm{Pic}_{X/S})_n$ est \emph{non ramifié}
\struck{\emph{étale}} sur $S$, et que \struck{$A_n$} $A_n$ est
\struck{ét.} un \emph{revêtement étale} de $S$. Donc si $p = 1$, i.e.\
\uncertain{$k(v)$} est de car.\ $0$, on a fini. Sinon, il faut regarder de
plus près - - - -
\note{« non ramifié », souligné, est écrit au-dessus de « étale », biffé.}

\marginal{?} À étudier : noyaux, images d'homomorphismes de schémas abéliens et
multiabéliens, en car.\ \uncertain{$\neq 0$}, notamment sur une
base non réduite…
\note{ce paragraphe est séparé du précédent par un trait horizontal sur toute
la largeur ; le point d'interrogation est dans la marge gauche. Le signe
après « car. » surcharge une lettre ; la fin de la dernière ligne sort du
feuillet.}

\section{Schémas abéliens sur un $S$ \emph{régulier}}
\note{titre de sa main, en tête du feuillet 22.}

\page{22}
a) Un schéma \uncertain{abélien} $P$ sur $S$ est projectif sur $S$
(Raisonnement de Nakai-\ill{})
\note{« abélien » est traversé d'un grand trait vertical qui en surcharge
l'initiale. Le second nom commence par une majuscule bouclée ; « Shimura »,
qui s'accorderait avec le titre du dossier, n'est pas exclu, mais le tracé ne
le porte pas assez pour qu'on l'imprime.}

b) Il existe un schéma de \uncertain{translations}, noté $G$, et $P$ est un
\uncertain{espace} \struck{fibré} principal homogène \struck{\ill{}}
isotrivial sous $G$.

c) Si $G$ est un schéma abélien sur $S$, alors
\[
  H^1(S, G) \subset H^1(K, G).
\]
Si $P_K$ est un fibré principal homogène sous $G_K$, pour qu'il
\struck{soit} provienne d'un $P$ pr.\ hom.\ sous $G$, il f.\ et s.\ qu'il
« se réduise bien ». En tous cas, il y a un plus grand \struck{\ill{}}
\uncertain{ouvert} $U$ \struck{\ill{}} dense de $S$ tel que $P_K$ provienne
d'un $P$ sur $U$ (i.e.\ l'ouvert où $P_K$ \struck{est} se réduit bien).
\marginal{\ill{}}
\note{le mot lu « ouvert » a la forme d'un trait discontinu terminé en
flèche, qui mène du mot noirci à $U$ ; dans la marge gauche, à hauteur de
cette ligne, un petit signe de renvoi.}

d) \marginal{résulte du th.\ de pureté} $S - U = Z$ est de codim.\ $1$ en
\struck{\ill{}} \ill{}.
\note{la note marginale est écrite en oblique à la hauteur de la fin de c) et
reliée par un trait à d).}

Donc
\[
  H^1(K, G) / H^1(S, G) \;\subset\;
  \sum_{\substack{D\ \uncertain{\text{div.}}\ \struck{\uncertain{\text{irréductible}}} \\
  \uncertain{\text{premier}}\text{ de } S}}
  H^1(\uncertain{K_D}, G) / H^1(\mathcal{O}_D, G)
\]
\note{le terme de droite est entouré ; de l'ovale descend une flèche vers
une colonne de mots soulignés dans le coin inférieur droit :
« \struck{\ill{}} : \uncertain{étudier} / \ill{} \ill{} / \ill{} \ill{} /
\ill{} \uncertain{complétion} ». L'indice $\mathcal{O}_D$ est souligné sur la
page.}

Analogue [\uncertain{th.\ de pureté}] pour la bonne réduction des schémas
abéliens ?
\note{cette question est prise dans une grande boucle tracée depuis « Donc » ;
« th. de pureté » est écrit au-dessus, sous une accolade qui le fait porter
sur « Analogue ».}

\end{document}
